\chapter{The Rationality Debate in the Psychology of Reasoning}
\label{sec:p13-rationality-1-2}
\index{rationality debate}\index{normative benchmark}\index{task interpretation}

Jonathan St. B. T. Evans's 2021 historical review asks how a field that tested
untrained people against formal logic came to investigate belief-based,
uncertain and context-sensitive reasoning.  Its answer is a history of
disagreement over the norm, the task understood by the participant, and the
reach of laboratory findings.  Evans is a participant in this history:
he defends a descriptivist alternative and identifies objections to his own
position.  His retrospective judgments about influence, a crisis or a new
paradigm are attributed interpretations, not a measurement of current consensus
\parencite[pp.~87--99]{Evans2021RationalityDebate}.

The longer conceptual history in Chapter~\ref{sec:p13-rationality-1-1}
has a different source and question.  Sturm's distinctions among
psychologism, description, prescription and justification are not merged
with Evans's particular history of twentieth-century reasoning experiments.
Nor does the chronology of logical discoveries in
Chapter~\ref{sec:p13-chronology-and-biography} supply a history of how people
responded to those calculi.

\section{What counted as rational reasoning}

Evans distinguishes instrumental rationality, serving an organism's or
person's goals, from epistemic rationality, acquiring and maintaining true
beliefs.  He explains why reliable knowledge can serve practical goals and
why truth-preserving deduction seemed a natural benchmark.  That rationale
does not make success at one experimental task a definition of either
rationality.  A valid argument preserves truth if its premises are true;
it does not establish that the premises are true or that making the inference
is the best use of limited resources.

The early studies cited in the review include Wilkins (1928), on changed
syllogistic content, and Woodworth and Sells (1935), on atmosphere effects.
Later reasoning experiments commonly used university students without formal
logic training, instructed them to assume supplied premises and judged whether
their conclusions necessarily followed.  The rationale was to investigate
inherent logical competence, not the effects of teaching a calculus.
Inhelder and Piaget's 1958 account made formal operations both a developmental
claim and a proposed description of adult thinking.  Wason in the 1960s
accepted logic as the norm but interpreted task failures as evidence of
illogicality and irrationality.  Those are distinct descriptive claims under
a partly shared benchmark \parencite[pp.~87--90]{Evans2021RationalityDebate}.

\begin{hbexample}{def:p13-r12-syllogisms}{Atmosphere, figure and a countermodel}
The review's first abstract example is:
some \(A\) are \(B\); some \(C\) are \(B\); therefore some \(A\)
are \(C\).  On two objects \(u,v\), take
\(A=\{u\}\), \(C=\{v\}\), \(B=\{u,v\}\).
Both premises are true and the conclusion false.  The universal conclusion
"all \(A\) are \(C\)" is false in the same model.
An existential atmosphere in the wording cannot supply the missing shared
witness.  The review's figural variant, some \(A\) are \(B\); some
\(B\) are \(C\); therefore some \(A\) are \(C\), has the same
countermodel.  Changing the arrangement of terms may change endorsement
behavior without changing this invalidity.

The two content examples from the review use these term assignments:
\begin{enumerate}
 \item \(A\): addictive things; \(B\): inexpensive things; \(C\): cigarettes.
 \item \(A\): millionaires; \(B\): hard workers; \(C\): rich people.
\end{enumerate}
Both have the form: no \(A\) are \(B\); some \(C\) are \(B\);
therefore some \(A\) are not \(C\).  Take
\(A=\{u\}\), \(B=\{v\}\), \(C=\{u,v\}\).
The premises hold and every \(A\) is a \(C\), refuting the conclusion.
All three terms are nonempty, so adding existential import for the universal
premise does not rescue the inference.
\end{hbexample}

Evans's review attributes endorsement figures of 71\% and 10\% for its two
content examples to Evans, Barston and Pollard's 1983 belief-bias study.
The chapter does not reproduce the full recruitment, denominator or uncertainty
record for those two percentages.  They are historical reported task figures,
not estimates of a population-wide inability to reason and not a new reanalysis
here.  Their argumentative role is to contrast the same logical form under
believable and unbelievable conclusions.  Logical validity, response criterion,
prior belief and comprehension must still be separated
\parencite[p.~88]{Evans2021RationalityDebate}.

Evans places mental-rule theories, including Braine--O'Brien and Rips, against
Johnson-Laird's mental-model approach, in which people imagine possibilities
and search for counterexamples.  The review describes mental-model theorists
as retaining much of propositional logic's normative role while explaining
failures through bounded working memory.  It contrasts this with researchers
who later changed the benchmark itself.  The full 2021 mental-model presentation
in Chapter~\ref{sec:p19-rationality-2-3} remains a distinct source-specific
account, not a theory reconstructed from these historical labels.

\section{Wason's tasks and the meaning of an error}

The 2--4--6 task, introduced by Wason in 1960, supplies a positive triple
and asks participants to discover the experimenter's rule through further
triples and yes/no feedback.  The intended rule is any ascending triple.
Participants can instead hypothesize equally spaced ascending triples and
continue to receive positive feedback from instances such as 1--3--5 or
10--30--50.  Evans notes the unusual procedure of recording reasons and
continuing after a wrong announced rule, and the original study's absence of
comparison conditions.  Rule discovery combines induction or abduction with
deductive testing; it is not simply another proof-validation item
\parencite[p.~89]{Evans2021RationalityDebate}.

For a concrete logical contrast, the triple 1--2--4 fits the broad ascending
rule but not the equally spaced hypothesis.  It would distinguish these two
hypotheses, whereas repeated equally spaced positive triples would not.
This does not show that all positive tests are useless.  Evans traces criticism
of Wason's verification-bias interpretation to Wetherick (1962) and later
Klayman and Ha (1987): the informativeness of a test depends on the relation
between the considered hypothesis and its alternatives.  The supplied
2--4--6 example biases that relation toward an unusually narrow hypothesis.
The historical dispute is about the chosen search norm and task ecology,
not about the logical fact that a counterinstance refutes a universal claim.

\begin{hbexample}{def:p13-r12-selection}{The four-card benchmark and its objective}
Wason's 1966 version describes cards with a letter on one face and a number
on the other, and asks whether every card follows the rule: if it has a vowel,
then it has an even number.  The visible faces are \(E,P,4,7\).
Let \(p\) be vowel and \(q\) be even.  Before any card is turned, the
vowel card could hide \(\neg q\), and the odd card could hide \(p\).
Those are the possible \(p\land\neg q\) violations.  A consonant is
not a violation whatever its number; an even number is not a violation
whatever its letter.  Thus choosing \(E\) and \(7\) is sufficient and,
with independent unknown reverse faces and this advance-choice objective,
neither can be omitted.  After finding one violation an adaptive search can
stop; that is a different decision stage.

The review reports low conventional solution rates, typically summarized
around 10\%, and common choices of the vowel alone or the vowel and even
number.  It does not supply a pooled-sample model making this number a
universal rate.  Choosing observations for expected information under a
probabilistic model, as in Oaksford and Chater's 1994 account, is a different
objective with priors, alternatives and observation assumptions.  It does
not make the material-rule violation criterion false.
\end{hbexample}

Wason linked both tasks to Popper's then-influential falsificationism and
argued strongly for irrationality.  Evans contrasts that interpretation with
Henle's 1962 view that apparent failures could reflect extra premises,
misinterpretations or rejection of the intended logical task.  The review
also qualifies thematic facilitation: not all realistic versions help, and
many successful variants change the question from whether a descriptive rule
is true to whether a permission or obligation has been obeyed.  An alcohol-age
rule is a deontic compliance problem; changing the content can change the
representation and objective as well as its familiarity
\parencite[pp.~89--90]{Evans2021RationalityDebate}.

The early 1970s heuristics-and-biases work introduced a parallel challenge in
judgment and decision making.  Evans contrasts it with the earlier intuitive-
statistician tradition associated with Peterson and Beach (1967) and the
influence of Edwards's subjective decision theory.  Representativeness,
availability and later prospect theory influenced reasoning researchers,
but Evans distinguishes Kahneman and Tversky's generally restrained claims
from stronger irrationality claims made by followers.  This distinction
prevents a psychological result from becoming a philosophical declaration
by attribution alone \parencite[pp.~90--91]{Evans2021RationalityDebate}.

\section{Cohen's critique: three different objections}
\label{sec:p13-r12-interpretation}

Cohen's 1981 question whether human irrationality could be experimentally
demonstrated received his negative philosophical answer.  Evans organizes
the objections into three categories, using his later 1993 reconstruction.
They attack different links in an inference from response to irrationality.

\begin{enumerate}
 \item \emph{Normative system.} Why does the chosen standard apply to this
       objective?  Popperian falsification, strict logical certification and
       Bayesian information acquisition are not identical norms.  Disputing
       one standard does not establish that whichever response is common
       must be correct.
 \item \emph{Interpretation.} Did participants represent the same premises,
       events and question as the experimenter?  If they add a premise or
       interpret an option differently, scoring against the intended problem
       can misdescribe their inference.  A plausible alternative interpretation
       still needs evidence; it is not an automatic explanation of every error.
 \item \emph{External validity.} Does the task and sample warrant the broader
       claim?  A deliberately deceptive item may be an unrepresentative
       problem, even when its recorded failures are real.  This objection
       does not show that every laboratory bias disappears outside the laboratory.
\end{enumerate}

The review's Linda example makes the second distinction explicit.  Let
\(T\) mean bank teller and \(F\) feminist, under one fixed information
state.  Probability requires \(P(T\land F)\leq P(T)\).
But if the bare teller option is understood as \(T\land\neg F\), the
two response options are not nested.  A distribution with masses
\(0.4,0.2,0.1,0.3\) on \(TF,T\neg F,\neg TF,\neg T\neg F\)
has \(P(TF)=0.4\leq P(T)=0.6\) while
\(P(TF)>P(T\neg F)=0.2\).  This demonstrates a possible change of
question, not that the original conjunction inequality can be violated.
The actual interpretation evidence and conflicting findings belong to
Chapters~\ref{sec:p19-conjunction-judgments} and
\ref{sec:p19-conjunction-language}; this toy distribution is not participant data.

Evans accepts the force of the three problems without accepting Cohen's
blanket conclusion.  He judges external validity the weakest general objection:
some demonstrations can mislead, but dismissing all bias results as contrived
laboratory tricks ignores their reported real-world and expert applications.
Conversely, the fact that a bias transfers somewhere does not establish that
every task effect transfers everywhere.  This is an attributed assessment in
the review, not an exhaustive new audit of all studies
\parencite[pp.~91--92]{Evans2021RationalityDebate}.

\section{The competing positions after 1981}

Stanovich's 1999 labels, as reconstructed by Evans, distinguish
\emph{Panglossians}, who deny a gap between normative and descriptive
rationality; \emph{Meliorists}, who allow a gap that education, capacity and
thinking dispositions can partly remedy; and \emph{Apologists}, who explain
departures through computational and other bounded limitations.  Baron and
Stanovich exemplify the meliorist position in the review.  These are historical
labels for arguments, not a personality scale or three exhaustive mechanisms.

Evans explicitly says he does not fit this division.  His descriptivist line,
from 1972 through collaborations with Over (1996, 2004) and Elqayam (2011),
questions the usefulness of normative assessment in explaining cognition.
Normative compliance may differ from instrumental success.  Stanovich
contests this distinction, treating normative theory as a description of
instrumental rationality.  Evans also records disagreement with his 2011
descriptivist proposal in its published commentaries: normativity retained
a central role even among researchers dissatisfied with classical logic
\parencite[p.~92]{Evans2021RationalityDebate}.

\subsection{Evolutionary mechanisms and ecological objectives}

Gigerenzer's early-1990s critique challenged narrow norms and explanations
in the heuristics-and-biases tradition; the review traces the exchange with
Kahneman and Tversky to 1996.  Frequency formats became a central example.
Cosmides and Tooby's 1996 evolutionary explanation competed with accounts
in which nested sets or simple mental models made the calculation easier.
The detailed 1995 frequency experiments are treated separately in
Chapter~\ref{sec:p19-natural-frequencies-bayesian-reasoning}: an improvement
from a format manipulation does not, by itself, choose its evolutionary
explanation.

In the fast-and-frugal programme described by Evans, Gigerenzer's 1996 and
1999 toolbox accounts portray low-cost heuristics as adaptive means of making
good decisions, rather than merely defective approximations to an unlimited
calculator.  His 2007 defense of gut feelings gives a stronger role to limited
information or ignorance and a weaker role to general reflective reasoning.
Evans emphasizes the resulting conflict with dual-process accounts, while
noting that different authors' uses of modules or heuristics need not share
every Fodorian criterion.  A successful heuristic in one environment is not
a theorem that ignorance or intuition is always preferable to further knowledge.

The review distinguishes Fodor's 1983 modularity from the later massive-
modularity proposal.  Fodor retained a general-purpose reasoning system;
Cosmides and Tooby proposed specialized evolved mechanisms and applied
cheater-detection accounts to social-contract versions of selection, later
including hazard avoidance.  Critics challenged both methods and the need
for that mechanism, proposing rivals without the evolutionary assumption.
The easier deontic task is not an isolated proof of a cheater-detection module.

Stanovich's 2004 objection separates the goals of individuals from reproductive
or ancestral success: an appetite useful in an earlier environment may be
unhelpful in the present, and contraception illustrates pursuit of individual
goals distinct from reproductive consequences.  Evans uses these as arguments
against conflating evolutionary and individual rationality, not as new medical
or population estimates.  Mercier and Sperber's 2011 proposal that reasoning
evolved for argumentation rather than simply problem solving occupies a
different position again; Evans notes a more complex development in their
2017 account.  It is not a synonym for either every adaptive explanation
or the entire fast-and-frugal programme.

Anderson's 1990--1991 rational analysis, developed for reasoning by Oaksford
and Chater, instead asks what goals and environment make a procedure optimal
at Marr's computational level.  In Evans's account, this treats environmental
adaptation as a starting assumption and lets a fitted objective challenge a
traditional logical norm.  A computational-level account is not yet a process
implementation, and descriptive fit does not establish the proposed objective's
normative authority.  Chapter~\ref{sec:p19-rationality-4-5} supplies Chater
and Oaksford's separate Bayesian presentation and its assumptions
\parencite[pp.~93--94]{Evans2021RationalityDebate}.

\subsection{Dual processes, dispositions and the limitations of a score}

Evans distinguishes automatic/intuitive Type~1 processing from reflective,
working-memory-dependent Type~2 processing, preferring those neutral labels
to names implying two unitary systems.  Type~1 may be learned, automated or
domain-specific.  Type~2 can fail through limited capacity, unsuitable
mindware or a failure to check an intuition.  Correctness is not the defining
feature of either type; intuition can be sound and deliberation mistaken.
Klauer's separate treatment in Chapter~\ref{sec:p19-rationality-2-5}
specifies alternative interaction architectures rather than inferring a
process from response speed alone.

Stanovich's individual-differences programme challenges a wholesale change
of norms by relating cognitive ability to conventional solutions, including
selection-task performance, while also emphasizing thinking dispositions.
In the review, IQ is not the same as rationality; the 2009 RQ proposal and
2016 CART project attempt a broader psychometric account.  Evans's 2021
assessment is that the CART still required item analysis and validation to
become a marketable instrument.  That is a dated judgment, not a statement
about any later test version or a new validation analysis here.

The chapter also records a challenge to the traditional explanation of
better performance by later reflection: some findings reviewed in De Neys's
2017 collection suggest rapid correct intuitions and individual differences
in those intuitions.  Evans considers some crisis language premature because
many such tasks are relatively simple and make few working-memory demands.
The disagreement therefore concerns interaction, task load and learning,
not a demonstrated equation between slow thought and rationality
\parencite[pp.~94--95]{Evans2021RationalityDebate}.

\section{What the new paradigm changed, and what it did not settle}
\label{sec:p13-r12-conditional-semantics}

Evans's 2002 critique named the deduction paradigm: untrained participants
assume given premises, judge necessary conclusions and are scored against
binary logic.  The review describes pressures from persistent content effects,
uncertainty and objections to treating natural-language conditionals as
material implication.  Its narrative contrasts retaining that benchmark and
accepting failures with changing the paradigm; the mental-model tradition
remained a significant resistance, so this was not unanimous abandonment.

In the source's Hillary-Clinton/UK--France illustration, the material
conditional \(p\to q\) is evaluated at a stipulated valuation in which
both \(p\) and \(q\) are false.  It is true because \(\neg p\lor q\)
is true.  This is not a current claim about officeholders or war, and it
does not refute the material connective's truth table.  It raises the
different question whether ordinary use of that sentence expresses that
connective.

The Ramsey-test proposal supposes \(p\) and evaluates \(q\) under that
supposition.  The associated psychological Equation is
\(P(\text{if }p\text{ then }q)=P(q\mid p)\), with positive probability
of the antecedent for the ordinary ratio conditional.  It is not the
probability of the material event: a stipulated distribution with
\(P(p\land\neg q)=0.1\), \(P(\neg p\land\neg q)=0.9\), and
zero on the other two cells has \(P(p\to q)=0.9\) but
\(P(q\mid p)=0\).  The chapter's 2003--2007 study references support its
historical account of the Equation's uptake; they do not establish a theorem
identifying every natural-language conditional with one probability rule.

Within the new-paradigm work, participants may reason from uncertain premises,
give graded conclusions and use relevant background beliefs.  Some researchers
emphasize pragmatics; others pursue probability, multivalued logic or Bayesian
decision norms.  Elqayam and Evans's 2013 strict/soft Bayesian distinction
separates normative commitments from more limited probabilistic descriptive
models.  Evans acknowledges that the new paradigm has no common settled
priorities, methods and theoretical frame.  Changing the benchmark is not
necessarily rejecting normativity altogether
\parencite[p.~95]{Evans2021RationalityDebate}.

\section{The review's concluding argument and a task lookup}

Evans ends with a methodological preference: describe cognitive and neural
mechanisms as elsewhere in cognitive psychology rather than repeatedly trying
to declare people rational or irrational.  He observes that memory limits,
visual illusions or inability to calculate large square roots are not normally
treated as global irrationality verdicts.  He also acknowledges that his
descriptivism was not the prevailing approach in the history he reconstructs.
The continued search for alternative formal norms in the new paradigm shows
that the normative question did not simply disappear.

His last evolutionary argument concerns the difference between what worked
in the past and what will meet future goals.  He contrasts human scientific
and cultural achievement with threats involving war, disease, climate and
weapons.  Those remarks are a philosophical warning about adaptation and
future success, not estimated disaster probabilities or a study proving that
evolution guarantees, or precludes, present rationality
\parencite[p.~96]{Evans2021RationalityDebate}.

\begin{hbmethod}{def:p13-r12-task-lookup}{Locate the disputed link before judging a response}
For a reasoning error claim, record the dated task, sample and instruction,
the intended formalization, and the observed response.  Ask whether the
dispute concerns the applicable norm, the interpretation of premises or
options, the implemented process, or external transfer.  A countermodel
settles an invalidity claim under one formalization; it does not identify
the cause of a participant's answer.  An alternative reading can change
the target comparison, but requires evidence that this was the reading used.

For an adaptive explanation, specify whose goal, the environment and the
optimization or satisficing criterion.  For a dual-process explanation,
specify the proposed functions and interaction, not merely quick versus
slow answers.  For a historical claim, distinguish the author's retrospective
assessment from an original position and from current consensus.  The
empirical discipline is the evidence record in
Definition~\ref{def:p19-evidence-record}; none of these checks is a permission
to select whatever norm makes the observed answer look correct.
\end{hbmethod}
